\documentclass[11pt]{article}

% ---------------------------------------------------------------
% arXiv-style preprint
% ---------------------------------------------------------------
\usepackage[margin=1in]{geometry}
\usepackage[T1]{fontenc}
% v26.3: pentagon glyphs for the Fig. 1 decoded-output marking, from
% DejaVuSans.ttf -- the figure's own font (local file, Bitstream-Vera
% licensed; newtx has no pentagon and this bundle's wasysym has only
% the hollow \pentagon with no filled variant). ORDER MATTERS: the
% family must be defined BEFORE newtxtext, which clobbers fontspec's
% file-based font lookup when loaded first (bisected, pent_test11).
\usepackage{fontspec}
\newfontfamily\DejaVuSansGlyphs{DejaVuSans.ttf}
\newcommand{\varpentagon}{\text{\DejaVuSansGlyphs\symbol{"2B1F}}}
\newcommand{\openpentagon}{\text{\DejaVuSansGlyphs\symbol{"2B20}}}
\usepackage{newtxtext,newtxmath}
\usepackage{graphicx}
\usepackage{booktabs}
\usepackage{array}
\usepackage{multirow}
\usepackage{amsmath}
\usepackage{xcolor}
\usepackage[font=small,labelfont=bf]{caption}
\usepackage{microtype}
\usepackage[hidelinks]{hyperref}
\usepackage{xurl}
\providecommand{\urlprefix}{}
\usepackage{enumitem}

\graphicspath{{../figures/}} % single source of truth: read figures/ directly, no copies in latex/figures (abolished 2026-10-22)

% float parameters: fig1 (image + long caption) exceeds the default
% \topfraction (0.7) of the text height, which blocked its [t] placement
% and queued every subsequent figure behind it, dumping all floats after
% the references. Allow tall top floats and fuller float pages.
\renewcommand{\topfraction}{0.95}
\renewcommand{\dbltopfraction}{0.95}
\renewcommand{\floatpagefraction}{0.85}
\renewcommand{\dblfloatpagefraction}{0.85}
\renewcommand{\textfraction}{0.05}

\newcommand{\wpm}{words/min}
\newcommand{\cpm}{chars/min}

\title{\bfseries Speech, Typing, and Handwriting
Neuroprostheses for People with Paralysis:\\
Data Openness, Independent Reproduction, and Clinical Maturity}

\author{we-linking}
\date{October 12, 2026}

\begin{document}
\maketitle

\begin{abstract}
\noindent
Severe paralysis from amyotrophic lateral sclerosis (ALS), brainstem
stroke, or spinal cord injury can leave a person fully alert yet unable
to speak or type, restricting communication with conventional assistive
devices such as eye trackers to a few words per minute. Implanted
communication neuroprostheses aim to restore natural-rate expression by
recording brain activity directly from the sensorimotor cortex while a
person attempts to speak, touch-type, or handwrite, and using artificial
intelligence decoders to translate those neural signals into text or
synthesized voice in real time. Here we review the 2021--2026 clinical
literature across all three modalities, supported by an open milestone
database and independent, from-scratch reproductions of all eight
publicly available datasets. We organize our assessment around three
questions:
\emph{(i)~Data openness:} Which studies openly share their human brain
recordings? Since 2023, studies using penetrating microelectrode arrays
have broadly released processed neural features on public repositories,
whereas studies using surface electrode grids (electrocorticography,
ECoG) and home-use recordings remain largely restricted, and no study
has released raw broadband voltages.
\emph{(ii)~Independent reproduction:} Can an outside laboratory
reproduce these landmark results on everyday hardware? Using a
modernized, lightweight decoding stack on a single 16\,GB consumer GPU,
we retrain and evaluate all eight open datasets from scratch, matching
or approaching published benchmarks on attempted speech, typing, and
handwriting and documenting larger gaps for inner speech. Session-level
analyses show when language modeling offsets neural decoding errors,
and memory profiling identifies what lets the pipelines run on a
consumer PC.
\emph{(iii)~Clinical maturity:} How many patients have actually
benefited, and how close is the field to a medical product? Across five
years of clinical trials and company demonstrations worldwide, only 20
people with paralysis have used a communication neuroprosthesis in real
time, and only one has used a system independently at home as a primary
communication device. Early systems collected thousands of training
trials from each participant. Recent studies have demonstrated daily
recalibration with tens of sentences and adaptation to a new user with
90 trials. Finally, we outline the clinical trial and
hardware durability requirements still needed for regulatory approval,
and explain how communication neuroprostheses must be redesigned to move
beyond English to tonal, character-based languages such as Mandarin
Chinese.
\end{abstract}

\section{Introduction and scope}
\label{sec:intro}

People who lose both limb movement and speech (anarthria or severe
dysarthria) after late-stage amyotrophic lateral sclerosis (ALS) or
brainstem stroke often retain normal cognition and language, yet can
communicate at only single-digit words per minute using conventional
augmentative and alternative communication (AAC) devices such as eye
trackers or head-tracking mice.
In the clinical study of Card et al., 2024, for example, participant
T15's residual speech delivered $6.8\pm5.6$ correct words per minute
(\wpm{}) and his head-mouse typing $6.3\pm1.3$ \wpm{}, compared with a
natural conversational pace of roughly 160 \wpm{} \cite{card2024}.
Implanted communication neuroprostheses bypass the paralyzed muscles by
recording neural activity directly from the brain's sensorimotor
cortex---the region that normally controls the vocal tract, hands, and
fingers---while the user attempts to speak, type on a keyboard, or write
letters by hand. Two classes of surgically implanted sensors are used in
current clinical trials: \emph{intracortical microelectrode arrays}
(such as the Utah array), which penetrate 1.5\,mm into the cortex to
record action potentials (spikes) from small populations of neurons, and
\emph{subdural electrocorticography (ECoG) grids}, which rest on the
surface of the brain beneath the dura to record pooled electrical
signals across a wider cortical area. Sequence-to-sequence neural
network decoders combined with language models then translate those
brain signals into on-screen text or synthesized speech.

Between 2021 and 2026, three communication modalities crossed the threshold
into \emph{online closed-loop} operation---meaning that brain signals
are decoded in real time so the user sees or hears the output as they
attempt to communicate, rather than having their recordings analyzed
\emph{offline} after the session ends: decoding attempted \emph{speech}
\cite{card2024, moses2021, metzger2023, willett2023, angrick2024, wairagkar2025, jude2026lis},
attempted \emph{touch-typing} \cite{jude2026typing}, and attempted
\emph{handwriting} \cite{willett2021, fan2023}. Two additional studies have
tested \emph{inner speech}---silently imagining words in one's head
without trying to move the mouth or vocal tract---as their primary task
\cite{wandelt2024, kunz2025}. Figure~\ref{fig:timeline} charts when and
at what performance level each milestone was reported. Despite dramatic
reductions in word error rate (WER, word substitutions, deletions, and
insertions divided by the reference word count) and daily setup time, major gaps remain between headline
demonstrations and routine medical care: clinical evidence still comes
from small single-center studies, public data sharing is uneven across
research groups, and almost all systems have been built exclusively for
English. This review evaluates the 2021--2026 literature around three
questions designed to be accessible across neuroscience, artificial
intelligence, and clinical medicine: First, which studies openly release
their neural recordings, and under what restrictions
(\S\ref{sec:openness})? Second, can an independent group reproduce the
published benchmarks from scratch on a single consumer-grade GPU, and
what limits performance and hardware feasibility
(\S\S\ref{sec:repro}--\ref{sec:feasibility})? Third, how many patients have benefited
in daily life, and what clinical and engineering milestones remain
before regulatory approval (\S\ref{sec:maturity})? We close with
the clinical and linguistic design choices for Mandarin speech
neuroprostheses (\S\ref{sec:chinese}).

\paragraph{Why this review includes only patients with paralysis.} We
include only studies enrolling participants with paralysis or severe
speech loss of neurological origin. Dozens of research groups worldwide
study speech decoding in able-bodied patients who have surface ECoG
grids or deep stereoelectroencephalography (sEEG) electrodes implanted
for a few days to localize seizures before epilepsy surgery or during
brain-tumor resection. By contrast, only a handful of clinical programs
worldwide implant chronic communication neuroprostheses in people with
paralysis. We exclude epilepsy and tumor monitoring studies because
their participants have intact speech and movement, their implants last
days rather than years, and combining those acute recordings with
chronic paralysis trials would distort every comparison of speed, error
rate, and device longevity. For the same reason, sEEG depth-electrode
studies are excluded outright: no clinical team implants penetrating
depth electrodes into deep brain structures in paralyzed patients to
build a communication device \cite{fogg2026}. Non-invasive recordings in
healthy volunteers (such as scalp electroencephalography [EEG] or
magnetoencephalography [MEG] \cite{levy2026}) are cited only as boundary
references. Finally, using brain signals to move a two-dimensional
computer cursor to click letters on an on-screen keyboard---from early
intracortical demonstrations \cite{hochberg2006} to peak speeds near 40
characters per minute (\cpm{}) \cite{pandarinath2017}---is a continuous
cursor-control problem rather than a language sequence-decoding problem,
and is included solely as a communication-speed baseline.

\paragraph{Relation to prior reviews.} Recent reviews have surveyed the
clinical milestones of speech neuroprostheses \cite{silva2024review},
the algorithmic evolution of neural speech decoding \cite{stavisky2025},
and the broader clinical landscape of invasive communication interfaces
\cite{khan2025}. Here we complement those syntheses in four ways:
(i)~a study-by-study audit of neural data openness (open, available on
request, or closed); (ii)~matched-scope independent reproductions of all
eight public datasets on a single consumer GPU, with transparent
accounting of where reproductions succeed or fall short; (iii)~a strict
patient-level maturity ledger that separates \emph{real-time closed-loop
participants} (twenty, including two company demonstrations) from
\emph{independent home users} (one); and (iv)~a unified comparison of
speech, touch-typing, and handwriting as three applications of a single
decoding pipeline borrowed from automatic speech recognition (ASR).

\section{Clinical milestones in speech, typing, and handwriting}
\label{sec:modalities}

Figure~\ref{fig:timeline} places speech, typing, and handwriting on a
common timeline. It records the tasks participants were instructed to
perform, the type of output produced, and whether they received
real-time feedback. These distinctions matter when interpreting the
reported vocabulary, speed, and error rates.

\begin{figure}[t]
\centering
\includegraphics[width=\linewidth]{main_fig_1.pdf}
\caption{\textbf{Timeline of speech, typing, and handwriting
neuroprostheses for people with paralysis, 2021--2026.} Markers sit at
each entry's publication, preprint, or announcement date along the
speech, typing, and handwriting tracks (open-data markers on the track
line; others offset toward their cards). Marker color (neural-data
availability; \S\ref{sec:openness}), fill (online closed-loop
vs.\ offline decoding), shape (electrode type), and card symbols for
speech ability ($\blacktriangle$/$\triangle$, after Silva et al., 2024a
\cite{silva2024review}), task instruction ($\blacklozenge$/$\lozenge$;
\S\ref{sec:defs}), and decoded output ($\varpentagon$/$\openpentagon$;
both on Metzger et al., 2023) follow the top legend. A $\star$ prefixes
the two studies meeting every term of the starred criterion---open
data, online closed-loop decoding, vocabulary $>$ 71.4k word forms,
online WER $<$ 12.8\%, and vocal-paralysis participants (Card et al.,
2024 \cite{card2024}; Jude et al., 2026a \cite{jude2026typing})---keyed
at the figure's bottom left. Each card states
its leading descriptor, speed, error rate (\texttt{n/r}: not reported),
vocabulary size, and participant IDs ($n$ counts all participants whose
recordings the study used; right-end totals count unique participants
per track). Section~\ref{sec:intro} gives the inclusion criteria.
Section~\ref{sec:cohorts} explains the counts and plotted exceptions.}
\label{fig:timeline}
\end{figure}

\begin{table}[t]
\centering
\footnotesize
\setlength{\tabcolsep}{3.6pt}
\renewcommand{\arraystretch}{1.12}
\caption{\textbf{The twenty participants of the
communication-neuroprosthesis literature and company reports.}
\textbf{Ch.\ (impl.)} = total implanted channels;
\textbf{Ch.\ (dec.)} = channels fed into the communication decoder;
\textbf{Interval} = symptom onset or diagnosis to implantation;
\textbf{Home use} = independent daily use.}
\label{tab:cohort}
\begin{tabular}{@{}>{\raggedright\arraybackslash}p{0.045\textwidth}>{\raggedright\arraybackslash}p{0.102\textwidth}>{\raggedright\arraybackslash}p{0.060\textwidth}>{\raggedright\arraybackslash}p{0.075\textwidth}>{\raggedright\arraybackslash}p{0.062\textwidth}>{\raggedright\arraybackslash}p{0.065\textwidth}>{\raggedright\arraybackslash}p{0.04\textwidth}>{\raggedright\arraybackslash}p{0.092\textwidth}>{\raggedright\arraybackslash}p{0.058\textwidth}>{\raggedright\arraybackslash}p{0.078\textwidth}>{\raggedright\arraybackslash}p{0.088\textwidth}@{}}
\toprule
ID & Institute & Etiology & Electrode & Ch.\newline (impl.) & Ch.\newline (dec.) & Side & Modality & Interval & Implanted & Home use \\
\midrule
T5 & BrainGate2 & SCI & Utah & 192 & 192 & L & handwriting & 9 y & 2016-08 & n/a \\
T12 & BrainGate2 & ALS & Utah & 256 & 128\newline 128 & L & speech\newline handwriting & 9 y & 2022-03 & Not yet \\
T15 & BrainGate2 & ALS & Utah & 256 & 256\newline 256 & L & speech & 5 y & 2023-07 & 3{,}800 h, 444 d \cite{card2026} \\
T16 & BrainGate2 & stroke & Utah & 256 & 64 & L & speech & 19 y & 2023-12 & Promising \cite{fogg2026} \\
T17 & BrainGate2 & ALS & Utah & 384 & 256\newline 256 & L & speech\newline typing & 3 y & 2024-02 & Promising \cite{jude2026typing} \\
T18 & BrainGate2 & SCI & Utah & 384 & 384\newline 384 & bi & typing\newline handwriting & 1.5 y & \mbox{2024-10}$^{\dagger}$ & n/a \\
T21 & BrainGate2 & ALS & Utah & 384 & 384 & bi & speech & 3 y & $<$2026-07 & Promising \cite{fogg2026} \\
T22 & BrainGate2 & stroke & Utah & 384 & 256 & L & speech & 3 y & $<$2026-07 & Promising \cite{fogg2026} \\
B1 & UCSF & stroke & ECoG & 128\newline 253 & 128\newline 253 & L & speech\newline speech & 16 y\newline 21 y & 2019-02\newline 2024-04 & Not yet \\
B3 & UCSF & stroke & ECoG & 253 & 253 & L & speech & 17 y & 2022-09 & Not yet \\
B6 & UCSF & ALS & ECoG & 253 & 253\newline 253 & L & speech & 4 y & 2025-03 & Not yet \\
CC01 & JHU & ALS & ECoG & 128 & 64 & L & speech & 8 y & 2022-07 & Not yet \\
P1 & Caltech & SCI & Utah & 288 & 96 & L & speech & 1.5 y & 2016-11 & n/a \\
P2 & Caltech & SCI & Utah & 320 & 64 & L & speech & 3 y & 2022-10 & n/a \\
Z1 & ZJU & SCI & Utah & 192 & 96 & L & handwriting & $\sim$3 y & 2019-08 & n/a \\
V1 & Neuralink & ALS & N1 & 1{,}020 & \emph{n/r} & L & speech & \emph{n/r} & $<$2026-10 & Not yet \\
V2 & Neuralink & ALS & N1 & 1{,}020 & \emph{n/r} & L & speech & \emph{n/r} & $<$2026-10 & Promising \cite{pouratian2026} \\
V3 & Neuralink & ALS & N1 & 3{,}072 & \emph{n/r} & L & speech & \emph{n/r} & $<$2026-10 & Promising \cite{pouratian2026} \\
P7 & Neuralink & ALS & N1 & \emph{n/r} & \emph{n/r} & \emph{n/r} & typing & \emph{n/r} & \mbox{$\sim$2025-06$^{\S}$} & n/a \\
C1 & Paradromics & MND & Connexus & \emph{n/r} & \emph{n/r} & \emph{n/r} & speech & \emph{n/r} & 2026-06$^{\ddagger}$ & Not yet \\
\bottomrule
\end{tabular}

\vspace{2pt}
Z1, V1--V3, C1 coined IDs $\cdot$ $^{\dagger}$derived from post-implant day $\cdot$ $^{\ddagger}$company report $\cdot$ $^{\S}$estimated from a \href{https://x.com/PairedWith_P7/status/2003575695663059181}{six-month follow-up post} $\cdot$ \emph{n/r} = not reported $\cdot$ n/a = not applicable $\cdot$ Not yet = performance below the Promising criteria $\cdot$ Promising = online, vocab $>$ 71.4k, WER $<$ 12.8\%, vocal paralysis.
\end{table}

\subsection{Tasks and representative studies}
\label{sec:defs}

\paragraph{Speech.} The standard paradigm is an instructed-delay copy
task: a sentence appears on screen, and after a go cue the participant
\emph{attempts to speak} it. ``Attempt'' covers a spectrum that must be
checked against each paper's operational definition rather than its
label: (1)~overt attempted speech (vocalized, possibly severely
dysarthric); (2)~mouthing (articulating without phonation);
(3)~intended speech (initiating articulation that the paralyzed
periphery cannot execute) \cite{metzger2023}; and (4)~inner speech
(covert verbalization with no articulation intent) \cite{kunz2025}.
Fig.~\ref{fig:timeline} separates this last, non-articulatory category
(filled $\blacklozenge$) from every articulation-engaging task (open
$\lozenge$); papers decoding under both instructions carry both
diamonds (Wandelt et al., 2024 \cite{wandelt2024}, internal and
vocalized speech; Kunz et al., 2025 \cite{kunz2025}, attempted and
inner speech; Jude et al., 2026b \cite{jude2026lis}, a three-phased
reading, internal-speech, and attempted-speech task).
On the vocabulary and speed axes, early closed-vocabulary systems
decoded 50 isolated words (Moses et al., 2021 \cite{moses2021}) or
1{,}152 words via silent NATO-codeword spelling (Metzger et al., 2022
\cite{metzger2022}). In 2023, open-vocabulary decoding arrived on the
intracortical line with the 125{,}000-word decoder of Willett et al.,
2023 (62 \wpm{} at 23.8\% WER \cite{willett2023}), while the ECoG line
demonstrated a 1{,}024-word talking avatar with synthesized voice
(Metzger et al., 2023; 78 \wpm{} at median 25.5\% WER
\cite{metzger2023}). Rapid daily recalibration then brought the
intracortical stack to 31.6 \wpm{} at a 9.8\% session-2 WER (2.5\% in
late sessions; Card et al., 2024 \cite{card2024}). Most recently, the
six-participant multi-user Transformer of Fogg et al., 2026 cut online
closed-loop WER to 3.9--8.0\% across three new users (a 55.9\% average
relative reduction over single-user models), rescued the locked-in user
T17 from a single-user 65.6\% WER (95\% CI 63.9--67.4) to 23.1\% (95\%
CI 21.5--24.7), and reduced new-user calibration to 90--200 sentences
\cite{fogg2026}, while Yoon et al., 2026 kept a deep-ensemble stack
running in real time on newly collected bilateral six-array recordings
from T12 (post-implant day 1{,}185 onward) \cite{yoon2026}.

On the voice-output track, Angrick et al., 2024 synthesized a six-word
vocabulary online from the Johns Hopkins CortiCom grid in a man with ALS
(naive listeners recognizing 80\% of the words with his voice profile
preserved \cite{angrick2024}), and Littlejohn et al., 2025 streamed a
personalized voice synthesizer continuously from Bravo-3's grid
(recurrent neural-network transducer [RNN-T] decoding in 80-ms steps,
synthesizing the participant's pre-injury voice at 58.8\% median
listener WER and 31.9\% text WER over a 1{,}024-word vocabulary
\cite{littlejohn2025}). On the intracortical side, Wairagkar et al.,
2025 achieved instantaneous synthesis in T15 across 8{,}489
cued-sentence trials. Naive listeners transcribed the synthesized audio
with a median WER of 43.75\% \cite{wairagkar2025}. An offline evaluation
of Brain2voice 2.0 on held-out trials from those recordings yielded a
mean listener WER of 5.24\% \cite{wairagkar2026}. Extending beyond voice
alone, Brosler et al., 2026 drove a personalized avatar with
simultaneous real-time attempted speech and covert (imagined) facial
gesture decoding from a single high-density grid per participant (online
simultaneous copy-task accuracies 70.0\% speech / 66.0\% gesture in
Bravo-6, and conversational blocks up to a median 100.0\% in Bravo-1r)
\cite{brosler2026}. Separately, Hou et al., 2025 showed that speech
motor cortex encodes decodable error signals during prosthesis use
(self-error decoding up to 86\%) \cite{hou2025}.

Beyond overt attempts, inner speech decodes online at 26--54\% WER (Kunz
et al., 2025 \cite{kunz2025}), and intended speech in a person with
locked-in syndrome reaches conversational WER $\sim$50\% over a
125{,}000-word vocabulary at 45--60 \wpm{} (Jude et al., 2026b
\cite{jude2026lis}). The first commercially manufactured entry, the
fully implanted wireless N1 of the Neuralink VOICE first-in-human trial
(NCT07224256), reported three adults with ALS decoding attempted speech
online: on a structured 125{,}000-word evaluation, the second and third
participants reached 54 \wpm{} at 6.6\% WER and 75 \wpm{} at 5.8\% WER
(the first participant, on the earliest electrode design, reached 45
\wpm{} at 7\% WER over a 220-word set) \cite{pouratian2026}. Two
mechanistic preprints from the BrainGate program map the planning code
these decoders read out: preparatory activity encodes diverse features
of intended movement \cite{rigotti2025}, and premotor cortex composes
phoneme-level representations into whole-word plans
\cite{abramovich2026}. Across the six BrainGate speech participants,
per-array ablations and 20-array whole-body mapping \cite{deo2026}
converge on the same anatomical gate (Fig.~\ref{fig:clinical}a):
speech decodes from ventral precentral area~6v (plus middle precentral
area~55b and primary motor area~BA4 in T15, T17, T21, and T22, including
T21's two right-hemisphere 6v/55b arrays, which alone decode at 48.0\%
offline multi-user WER), whereas Broca's area~44 (T12, and excluded a
priori in T22), the premotor eye field (T16), and dorsal hand-knob
area~6d (T16, T17, and T22) stay at or near chance for speech
\cite{willett2023, jude2026lis, kunz2025, fogg2026, deo2026}.

\paragraph{Typing.} The participant attempts to touch-type on a (real or
on-screen) QWERTY keyboard; the decoded object is the keystroke sequence
itself, not a cursor and not dictated speech. Jude et al., 2026a (typing)
is the anchor:
T18 reached 110 \cpm{} (22 \wpm{}) at 1.6\% WER---81\% of the
age-matched able-bodied smartphone-typing rate of 27.3 \wpm{}
\cite{palin2019}---with 30 calibration
sentences, drifting only to 12\% WER after a week without recalibration
\cite{jude2026typing}. Its technical ancestry (bimanual RNN control,
finger decoding, premotor geometry; BRAND real-time platform) is
documented in the typing-neuroprosthesis literature
\cite{deo2024, willsey2025, shah2025, ali2024}; a preprint from the same
program traces how typing skill consolidated in dorsal premotor cortex
over practice \cite{hashimoto2026}. A Neuralink company update shows
P7 (``Jake'', its seventh implant recipient, who has ALS and retains
speech) typing online by imagining movements of all ten fingers. The
company describes decoding signals for both hands from a unilateral N1
implant and says Jake uses the system for everyday computer tasks.
Neuralink reports speeds of up to 40 \wpm{} among participants and
provides no WER \cite{neuralink2026, schneider2025}.

\paragraph{Handwriting.} The participant attempts to handwrite
characters; BrainGate's T5 (C4 spinal cord injury) is the only
participant ever to run this modality online. Willett et al., 2021: 90
\cpm{} at
94.1\% raw online accuracy (5.9\% character error rate, CER; 25.1\% raw
online WER), 0.89\% CER (3.4\% WER)
after offline language-model correction, and 73.8 \cpm{} free-text copy
at 8.54\% CER---against an age-matched smartphone-typing reference of
115 \cpm{} \cite{willett2021}. Fan et al., 2023 demonstrated plug-and-play
stability across 403 days of monthly evaluation sessions (15 in total;
73.2 $\pm$ 11.0 \cpm{}, 6.16\% online WER over
a 130{,}000-word vocabulary) \cite{fan2023}.
Notably, the 2021 handwriting architecture has since been absorbed by
the speech stack: Levin et al., 2026 retrained handwriting data with the
GRU+CTC speech pipeline and reached 2\% CER in a single session (and
23--40\% speech PER with cross-brain transfer at 50--84 target-user
sentences, evaluated without a language model)
\cite{levin2026}. Cross-participant generalization is also the explicit
target of second-hand reanalyses---cross-subject pretraining
\cite{boccato2026} and adversarial session alignment \cite{align2026}---
which point the same direction at lower cost
(Table~\ref{tab:crosssubject}). A parallel Chinese
trajectory-decoding line
(stroke-level trajectories rather than character sequences, with
video-guided per-stroke supervision) started offline and has now
reached online closed-loop operation: Xu et al., 2025
reach 91.1\% closed-set recognition over a 1{,}000-character template
library by multi-day-fused DTW matching (plotted on the handwriting
lane of Fig.~\ref{fig:timeline}), and the follow-up Yang et al., 2026
closes the loop on the same participant---an LSTM decoder writes
trajectories in real time and a handwriting-BERT classifier, pre-trained
on 6{,}825 character classes from 122 human writers (HIT-OR3C), reports
94.45\% online accuracy over a 295-character vocabulary at 8.55 $\pm$
2.77 \cpm{} (plotted on the same lane) \cite{xu2025, yang2026}; the
companion study of Wang et
al., 2026 maps the cortical representation of multidimensional
handwriting movement \cite{wang2026}.

\subsection{Clinical programs and participant counts}
\label{sec:cohorts}

Table~\ref{tab:cohort} lists all twenty individuals with paralysis who
participated in the studies or company demonstrations covered through
2026. Eighteen participants appear in peer-reviewed papers or preprints.
The studies span several clinical trial programs, including the multi-site BrainGate2 intracortical consortium (T5, T12,
T15, T16, T17, T18, T21, and T22), the University of California, San
Francisco (UCSF) ECoG program (B1, B3, and B6; the bilingual study of
Silva et al., 2024b \cite{silva2024} and the high-density re-implant of
Brosler et al., 2026 \cite{brosler2026} both involve B1, first reported
by Moses et al., 2021 \cite{moses2021}), the Johns Hopkins CortiCom ECoG
trial (CC01 \cite{angrick2024, luo2023}), the Caltech supramarginal gyrus
intracortical trial (P1 and P2 \cite{wandelt2024}), the Zhejiang
University Chinese handwriting trial (Z1 \cite{xu2025, yang2026}), and the
wireless N1 intracortical implant of the Neuralink VOICE clinical trial
(V1--V3 \cite{pouratian2026}). The remaining two entries are
company-reported demonstrations without peer-reviewed papers or published
error rates: Neuralink's seventh implant recipient, P7 (``Jake''), who
performs imagined ten-finger typing \cite{neuralink2026, schneider2025},
and Paradromics' first Connect-One participant, C1, who decodes
attempted speech in real time using the Connexus cortical module
\cite{perkins2025, paradromics2026}. All twenty received their implants
under registered clinical trials (seven on ClinicalTrials.gov and one in
the China Clinical Trial Registry). Alongside these first-hand patient
studies and two preprints introducing new participant recordings
\cite{yoon2026, levin2026}, Figure~\ref{fig:timeline} also tracks nine
reanalyses of openly released datasets. Four date from 2024--2025
\cite{feng2024, falcon2024, li2025dcond, feghhi2025} and five from 2026
\cite{wairagkar2026, boccato2026, align2026, bit2026, feghhi2026}.
The original study team carried out the Brain2voice 2.0 reanalysis.

The timeline counts participants once per modality and includes those
whose recordings were decoded offline. Of the twenty unique
participants, 16 contributed to speech, 3 to typing, and 4 to handwriting,
with T12, T17, and T18 contributing to two modalities. Six participants
retain natural speech---the five with spinal cord injury (SCI; T5, T18,
P1, P2, Z1) plus P7---and fourteen have vocal paralysis
\cite{silva2024review}.

Twenty unique participants qualified for the online closed-loop count:
sixteen in speech, three in typing, and two in handwriting (T5 and Z1),
with T17 counted in both speech and typing. This total includes two
company-reported demonstrations (P7 and C1)
\cite{neuralink2026, paradromics2026}. Several participants returned for
multiple studies over years, so a new publication can mark a technical
advance without adding a new participant. Evidence for independent home
use is considered in \S\ref{sec:home_use}.

\section{Public datasets and decoding benchmarks}

Figure~\ref{fig:vocab} starts with vocabulary coverage across the
field (a), then narrows to the public datasets and their reported online
performance (b--d). The final two panels use those releases to compare
independent reproductions with published benchmarks
(\S\ref{sec:repro}).

\begin{figure}[tp]
\centering
\includegraphics[width=\linewidth]{main_fig_2.pdf}
\captionsetup{singlelinecheck=false,font=footnotesize,skip=2pt}
\caption{\textbf{Our reanalysis of the open-sourced datasets: vocabulary
funnel to independent reproduction on one 16\,GB consumer GPU.}
(a) Maximum decoder vocabulary size (inflected forms; dashed:
open-ended output, no fixed lexicon; \cite{xu2025} omitted) against
native-speaker tiers \cite{brysbaert2016}. (b) Labeled sentences
across the 8 open-data studies, 11 splits. (c) Online WER from daily summed word edits divided by reference words,
averaged equally across days (error bars: sample SD; listener median
for voice synthesis). (d) Online
speed against able-bodied speech, typing, dwell
eye-typing, and point-and-click BCI references
\cite{card2024, willett2021, pandarinath2017, palin2019, dhakal2018, hamid2024}.
(e) Before language modeling, our offline reproduction (blue) versus
paper benchmarks (teal; gray tags indicate online or offline origin):
greedy CTC PER, raw CER, or $1-r$ mel gap for the Wairagkar voice arm.
(f) Our offline full-pipeline WER (blue) alongside the online WER
and error bars from (c) (teal). For Wairagkar et al., 2025, only the
online listener WER is available (Table~\ref{tab:repro}).
Values and SDs for (b)--(f): Source Data file.}
\label{fig:vocab}
\label{fig:repro}
\end{figure}

\subsection{Vocabulary coverage and data availability}
\label{sec:openness}

Across the studies in Fig.~\ref{fig:vocab}a, the decoded sets span 6--50
word forms in the early ECoG demonstrations
(Angrick et al., 2024; Wandelt et al., 2024; Moses et al., 2021) and
1{,}024 word forms in the 2023 avatar and 2025 voice work (Metzger et
al., 2023; Littlejohn et al., 2025), while the 125{,}000-word $n$-gram
language model of Willett et al., 2023 has been the workhorse of
open-vocabulary decoding ever since (ten studies in
Fig.~\ref{fig:timeline}). Four studies go further and decode open-ended
output with no fixed lexicon at all: the
handwriting systems of Willett et al., 2021
\cite{willett2021} and Levin et al., 2026 \cite{levin2026}, and the
voice synthesis of Wairagkar et al., 2025 (online) and 2026 (offline)
\cite{wairagkar2025, wairagkar2026} (Box~1 explains how machine word-form
lexicons compare against human vocabulary benchmarks).

\vspace{4pt}
\noindent\fbox{\begin{minipage}{\dimexpr\linewidth-2\fboxsep-2\fboxrule\relax}
\small\textbf{Box 1 $|$ Counting vocabulary size: word forms, lemmas,
and word families.}\\[2pt]
What a reported vocabulary size means depends on what counts as a word,
and the three conventions in play differ by multiples. Take the six
spellings \emph{walk, walks, walked, walking, walker, walkers}: a
machine lexicon counts every distinct spelling (word forms)---six
entries; a dictionary counts lemmas, one headword per meaning with its
inflections merged (\emph{walk} covers the first four, \emph{walker} the
last two)---two; and psycholinguistic vocabulary studies typically count
word families, one entry per root regardless of derivation---one. Every
fixed vocabulary in Fig.~\ref{fig:vocab} is a flat word-form list: the
shared 125{,}000-word language model literally lists \textsc{walk,
walks, walked, walking} as four separate entries \cite{jude2026typing}.
The like-for-like human anchor is thus the word-form tier of Brysbaert
et al.\ \cite{brysbaert2016}: an average 20-year-old native speaker of
American English knows 71{,}400 inflected word forms ($1.7\times$ their
42{,}000 lemmas, rooted in 11{,}100 word families)---so the
125{,}000-word decoder lexicon already exceeds a native speaker's
repertoire under the same counting unit.
\end{minipage}}
\vspace{4pt}

Figures~\ref{fig:timeline} and~\ref{fig:vocab}a distinguish public
neural recordings, recordings available on request, and closed data.
Eight public releases underpin the reproduction benchmarks in
\S\ref{sec:repro}. First, \emph{the intracortical line went
broadly open from 2023 onward}: Willett et al., 2023 (Brain-to-Text~'24
benchmark; Dryad, $\sim$37\,GB), Fan et al., 2023 (handwriting
plug-and-play; Dryad, 3.57\,GB), Card et al., 2024 (Brain-to-Text~'25;
Dryad, 11\,GB), Wairagkar et al., 2025 (26\,GB), Kunz et al., 2025,
Jude et al., 2026a (typing; 35.4\,GiB), and Jude et al., 2026b
(locked-in syndrome) are all directly downloadable
\cite{card2024, willett2023, wairagkar2025, jude2026lis, jude2026typing, fan2023, kunz2025}, alongside the earlier Willett et al., 2021
handwriting dataset (1.31\,GiB \cite{willett2021}), the Caltech
supramarginal gyrus internal-speech dataset of Wandelt et al., 2024 on
Zenodo \cite{wandelt2024}, and the one fully open ECoG dataset in our
cohort, Angrick et al., 2024 (Johns Hopkins subdural grids; OSF neural
data and audio) \cite{angrick2024}. The exception on the intracortical
side is the Chinese handwriting line: both Xu et al., 2025 and the
online follow-up Yang et al., 2026 gate their neural recordings behind
request procedures, citing clinical-trial protocol and participant
privacy \cite{xu2025, yang2026}. Second, \emph{the UCSF ECoG line and
home-use neural recordings remain restricted or closed}: Moses et al.,
2021, Metzger et al., 2022/2023, Silva et al., 2024b, Littlejohn et al.,
2025, and Brosler et al., 2026 gate their neural recordings behind
application procedures or restricted-access repositories (or release
only figure-level summary data)
\cite{moses2021, metzger2023, metzger2022, littlejohn2025, brosler2026, silva2024}, and the home-use neural recordings of Card
et al., 2026 are withheld for participant privacy, with summary CSVs
released \cite{card2026}. Third, \emph{every publicly released neural
dataset consists of derived features}---binned threshold-crossing counts
and spike-band power---never raw wideband voltages or spike-sorted
units; de-identification at the feature level means the pre-processing
stage cannot be audited post hoc.

``Open'' also comes in degrees that matter for reproduction: the
Brain-to-Text~'25 test labels remain unreleased at the time of writing;
Wairagkar et al., 2026 reuse the public Dryad recordings of Wairagkar et
al., 2025 \cite{wairagkar2026}; and the 2026 preprints have not yet
released their newly collected neural recordings---including the 20 new
sessions in Levin et al., 2026, the six-array T12 recordings in Yoon et
al., 2026, the six-user corpus in Fogg et al., 2026, and the recordings
in Nason-Tomaszewski et al., 2026
\cite{fogg2026, yoon2026, levin2026, nason2026} (whereas Hahn et al., 2025
released their dataset on Dryad upon publication \cite{hahn2025}).

\subsection{Released corpus sizes}
\label{sec:corpus_size}

Figure~\ref{fig:vocab}b counts labeled sentences in the eight public
releases used in our reproductions. Willett et al., 2021 released
1{,}000 sentences (8{,}198 word tokens) across 10 of 11 sessions.
Card et al., 2024 released 45 sessions comprising 9{,}498 transcribed
training or validation trials (60{,}091 word tokens) and 1{,}450
unlabeled test trials, for 10{,}948 trials in total. Jude et al., 2026b
released 3{,}674 sentences (20{,}237 word tokens) across 13 of 14
sessions. We casefolded tokens, removed punctuation, and counted only
trials with reference transcripts. This excludes 1{,}200 blind
hold-out trials in Brain-to-Text~'24 and 1{,}450 in Brain-to-Text~'25.

Amounts collected over a study and amounts released in a public
dataset describe different things. Table~2 of Fogg et al., 2026
quantifies the intracortical corpus per participant
(Supplementary Fig.~S2c): T12
10{,}880 copy-task trials; T15 19{,}557 copy trials \emph{plus} 260{,}424
self-generated sentences accumulated across independent home use
(278{,}531 training trials over 748 days); T15's public
Brain-to-Text~'25 subset 10{,}948 trials; T16 11{,}274, T17 3{,}763,
T21 2{,}942, and T22 4{,}508 \cite{fogg2026}. In task time, the copy
corpus of Card et al., 2024 alone spans 67.4\,h of attempted speech
(248.3\,h including self-generated conversation). For
reference, Moses et al., 2021 collected 22.5\,h of isolated words
(9{,}800 trials) plus 2.1\,h of sentences \cite{moses2021}, and Metzger
et al., 2023 about 1.6\,h/day for 14 days (9{,}506 trials)
\cite{metzger2023}.

The 10{,}850 sentences reported by Willett et al., 2023, for example,
correspond to an 8{,}800+880 split in the release cited here. Card et al.,
2024 reported 10{,}481 copy sentences across 50 sessions, while the
public release contains 10{,}948 trials across 45 sessions. Task hours
offer a more consistent measure of recording burden across studies
because isolated-word and sentence trials differ in duration.

\subsection{Online error and communication speed}

Panels c and d of Figure~\ref{fig:vocab} summarize day-level online
WER and communication speed. Handwriting and typing have reached
78--81\% of
age-matched able-bodied smartphone-typing speeds (90 vs.\ 115 \cpm{}
for T5 \cite{willett2021, palin2019}; 22 vs.\ 27.3 \wpm{} for T18
\cite{jude2026typing, palin2019}), and 35--61\% of the population-wide
anchors (36.2 \wpm{} smartphone, 51.6 \wpm{} keyboard
\cite{palin2019, dhakal2018}),
while open-vocabulary speech remains a factor 2--3 below conversational
pace but roughly an order of magnitude above the patients' own AAC
baselines. Error rates describe modality-specific tasks. For example,
the 1.6\% typing WER was measured with a 30-symbol output set, built-in
word segmentation, and templated test sentences \cite{jude2026typing}.

Raw speed counts output that users may need to correct.
Converting each headline result
to \emph{correct words per minute}, $\text{rate}\times(1-\text{error})$,
using each paper's own paired speed--error anchor (with character-rate
systems converted at the five-characters-per-word convention of
Jude et al., 2026a, typing \cite{jude2026typing}),
tightens the ranking: the best speech systems deliver 47--71 correct \wpm{}
(the fully implanted Neuralink system reaching 70.7 in a conference
abstract \cite{pouratian2026}; Metzger et al., 2023 reaching 58.1
\cite{metzger2023}; Card et al., 2026 home use reaching 48.0 silent and
$\gtrsim$30 vocalized \cite{card2026}; Willett et al., 2023 reaching
47.2 \cite{willett2023}; Card et al., 2024 session~2 reaching 28.5
within its 2.5--9.8\% online WER range \cite{card2024}; Moses et al.,
2021 reaching 11.3 \cite{moses2021}),
handwriting and typing deliver
13.7--21.6 (Jude et al., 2026a: 21.6 \cite{jude2026typing}; Willett et
al., 2021: 16.9 \cite{willett2021}; Fan et al., 2023: 13.7 from 73.2
\cpm{} at WER 6.2\%
\cite{fan2023}), and every BCI system still clears T15's own 6.3--6.8 correct
\wpm{} AAC baselines by $2\times$ to $11\times$---none reaches half
of able-bodied conversational pace, and among open-vocabulary systems only
the Neuralink row crosses the able-bodied keyboard anchor (51.6 \wpm{}).

\paragraph{Public benchmark competitions.} Two public releases also
supported competitive benchmarks.
Brain-to-Text~'24 ran on EvalAI with the
T12 data: a private 1{,}200-sentence hold-out scored by aggregate WER, a
declared baseline of 9.7\% (RNN CTC $\to$ 5-gram $\to$ LLM rescoring),
and a winning entry at 5.77\% that decoded diphones ($\sim$1{,}600
classes) and merged ten RNNs with a fine-tuned LLM; all top-three
entries relied on LLM-merged ensembles, Mamba and Transformer entries
did not beat the RNN baseline, and the benchmark itself remains open
\cite{btot24lessons}---its EvalAI board listed 49 teams as of
2026-10-16, with post-close submissions pushing the open-board best to
3.34\%, ahead of the closed competition's 5.77\% \cite{evalai24board}.
Its successor, Brain-to-Text~'25, moved to Kaggle
with the T15 data and cash prizes (2025-07-07 through 2026-01-01;
\$5{,}000/\$3{,}000/\$1{,}000 most-innovative, sponsored by Blackrock
Neurotech): \mbox{576 entrants} and 466~leaderboard teams competed on the
1{,}450-sentence test partition, against a declared baseline of 6.70\%
WER and a final private-leaderboard best of 1.54\% \cite{btot25kaggle}.
These are offline numbers, earned with unlimited compute and ensembling
on a static test set---Card et al., 2024's online session-2 anchor of
9.8\% remains the clinical reference \cite{card2024}---but as
standardized external yardsticks they are exactly what the post-2024
open-data turn was for. The hosts then closed the loop themselves:
Yoon et al., 2026---authored by the organizers (Stanford-NPTL +
UCD-NPL, with Emory/GaTech and Brown)---ran the boards' winning
recipe, LLM-merged deep ensembles, as the first real-time closed-loop
test on T12's 2025 bilateral six-array re-implant recordings, cutting
online WER from 33.7\% to 26.0\% at 65.4 wpm
\cite{yoon2026} (Fig.~\ref{fig:clinical}c).

\section{Independent reproduction of published benchmarks}
\label{sec:repro}

Figure~\ref{fig:vocab}e compares the neural stage before language
modeling, and panel f shows our offline end-to-end word error alongside
the online results in panel c. We retrained on each public release
and matched dataset splits, decoding stages, and evaluation statistics
to the extent the releases allowed. Table~\ref{tab:repro} gives the
online results and the offline reproduction values.

\subsection{Shared decoding pipeline}

Text decoders across the three modalities share a cascade adapted from
automatic speech recognition (ASR). On the sensor side, the intracortical line
froze its core front end in 2021: no spike sorting; threshold crossings in
20\,ms bins (handwriting \cite{willett2021, fan2023}, $1\times N_{\text{ch}}$)
concatenated from 2023 onward with spike-band power (speech and typing
\cite{card2024, willett2023, jude2026typing}, $2\times N_{\text{ch}}$),
Gaussian-smoothed, blockwise $z$-scored, with a
day-specific affine input layer absorbing cross-day nonstationarity. The ECoG line
instead uses high-gamma envelopes \cite{moses2021}. Decoding then walks
the ASR textbook: a stacked GRU (five layers of 512--768 units;
Transformer variants from 2026 \cite{fogg2026, card2026}) trained with
connectionist temporal classification (CTC) \cite{graves2006} emits
frame-level label probabilities; a beam search over a lexicon-constrained
$n$-gram language model turns them into words; and an n-best list is
rescored by a large language model. Early text decoders used cross-entropy classification
\cite{moses2021, willett2021, metzger2022, silva2024}. CTC became common in
text decoding from 2023 onward
\cite{card2024, metzger2023, willett2023, jude2026lis, jude2026typing, kunz2025, fogg2026, nason2026}.
For streaming voice synthesis, Littlejohn et al., 2025 used an
RNN-Transducer (RNN-T) \cite{littlejohn2025}.
Jude et al., 2026a adapted the speech stack for typing with a 30-key
output head, and Levin et al., 2026 used the stack for handwriting
\cite{levin2026}. By 2026 the text-decoding stacks shared much of their architecture,
with modality-specific front ends and output units (Box~2).

\vspace{4pt}
\noindent\fbox{\begin{minipage}{\dimexpr\linewidth-2\fboxsep-2\fboxrule\relax}
\small\textbf{Box 2 $|$ What communication neuroprostheses borrowed from automatic speech recognition.}\\[2pt]
The neuroprosthesis decoding cascade borrows almost everything from its
industrial source. The language-model decoding engine is not
merely ASR-inspired: it is a literal code fork of the \texttt{wenet}
toolkit \cite{wenet} (Mobvoi, 2020; Apache-2.0), whose C++ runtime,
weighted finite-state transducer (WFST) beam search, and bundled SRILM
scripts were carried through \texttt{handwritingBCI} (2021),
\texttt{speechBCI} (2023), and \texttt{nejm-brain-to-text} (2024) with
roughly 2{,}578 files shared and about nine modified---the C++ decoder
input was simply swapped from acoustic features to recurrent neural
network (RNN) phoneme probabilities. The shared text decoder uses
CTC, one of the two training paths in \texttt{wenet}'s U2 architecture.
Where production ASR trains on corpora such as WenetSpeech's 10{,}000+ hours
of multi-domain web speech \cite{wenetspeech}, the neuroprosthesis stack
trains on 22.5--67.4\,h from a single speaker, for whom per-day
recording nonstationarity, not data volume, dominates the engineering.
\end{minipage}}
\vspace{4pt}

\subsection{Reproduction protocol and engineering choices}

All of our reproductions run on a single Windows consumer PC (RTX 5060
Ti, 16\,GB VRAM; 64\,GB RAM) with WSL, directly on the public Dryad
releases (SHA-256 verified file by file) and the official code
repositories. Two engineering substitutions define the effort, both of
which the field itself made in 2026:

\emph{KenLM instead of Kaldi/WFST as the default decoding engine.} The
official $n$-gram stack is a \texttt{wenet} C++ fork whose composed
token--lexicon--grammar (TLG) graph for the Willett 5-gram spans
$\sim$44\,GB; direct loading exhausts a 20\,GB WSL allocation in 26
seconds, so anchor-reproduction runs use memory-mapped access to the
official graph (for scope comparability), while new analysis uses KenLM
+ Flashlight-Text---the same substitution the field itself made in 2026
(\S\ref{sec:memory}) \cite{fogg2026}.

\emph{Qwen-class instead of GPT/OPT rescorers.} OPT-6.7B in 16-bit
($\sim$13\,GB) barely fits a 16\,GB card; Qwen3.5-4B in bfloat16
(8--10\,GB) is the field's own replacement \cite{fogg2026}, and local
Qwen3-4B serving has replaced remote GPT-4.1-mini calls in the ensemble
literature \cite{yoon2026}.

A third, unplanned substitution became necessary: the official TF1
codebases of Willett et al., 2021 and of Jude et al., 2026a (typing)
offer no Windows GPU path, so we ported them to PyTorch; the Jude port
runs $47\times$ faster with data conversion byte-identical to the
official pipeline, and both ports have since been retrained from
scratch (Table~\ref{tab:repro}).
For scale, the complete 1{,}426-trial, 41-session evaluation of the
official Card checkpoint (greedy PER 10.19\%) re-runs end to end in
about 80 seconds on this GPU---acoustic inference is no longer the
bottleneck of the stack.

\begin{table}[t]
\centering
\footnotesize
\setlength{\tabcolsep}{3pt}
\renewcommand{\arraystretch}{1.12}
\caption{\textbf{Released data and decoding performance across eight
open-data datasets (11 participant splits).}
Online evaluation days follow Fig.~\ref{fig:vocab}c; reproduced WER is
offline after language modeling. Values are mean $\pm$ SD where applicable.}
\label{tab:repro}
\begin{tabular}{@{}>{\raggedright\arraybackslash}p{3.50cm}lrrrlll@{}}
\toprule
Study & Participant & \shortstack[r]{Released\\sentences} & \shortstack[r]{Total\\recording\\days} & \shortstack[r]{Online\\evaluation\\days} & \shortstack[l]{Online speed\\(words/min)} & \shortstack[l]{Online WER\\(\%)} & \shortstack[l]{Reproduced\\WER (\%)} \\
\midrule
Willett et al., 2021 \cite{willett2021} & T5 & 1{,}000 & 10 & 5 & $15.4{\pm}2.2$ & $23.9{\pm}8.4$ & $10.0{\pm}3.58$ \\
Willett et al., 2023 \cite{willett2023} & T12 & 9{,}680 & 24 & 5 & $65.3{\pm}0.9$ & $23.8{\pm}1.8$ & $23.58{\pm}5.25$ \\
Fan et al., 2023 \cite{fan2023} & T5 & 1{,}315 & 21 & 10 & $14.6{\pm}2.2$ & $4.0{\pm}1.2$ & $4.10{\pm}1.20$ \\
Card et al., 2024 \cite{card2024} & T15 & 9{,}498 & 45 & 14 & $32.5{\pm}2.7$ & $4.9{\pm}3.3$ & $6.25{\pm}6.54$ \\
Wairagkar et al., 2025 \cite{wairagkar2025} & T15 & 8{,}489 & 37 & 8 & --- & 43.8 & --- \\
Kunz et al., 2025 \cite{kunz2025} & T12 & 500 & 1 & 1 & --- & 22.9 & $94.35{\pm}0.0$ \\
Kunz et al., 2025 \cite{kunz2025} & T15 & 720 & 3 & 2 & --- & $43.3{\pm}23.8$ & $78.07{\pm}0.59$ \\
Kunz et al., 2025 \cite{kunz2025} & T16 & 450 & 2 & 1 & --- & 44.4 & $82.38{\pm}0.0$ \\
Jude et al., 2026a \cite{jude2026typing} & T17 & 668 & 8 & 6 & $6.8{\pm}1.8$ & $10.5{\pm}3.1$ & $11.41{\pm}6.85$ \\
Jude et al., 2026a \cite{jude2026typing} & T18 & 1{,}033 & 9 & 7 & $14.0{\pm}4.5$ & $3.5{\pm}0.6$ & $6.66{\pm}3.67$ \\
Jude et al., 2026b \cite{jude2026lis} & T17 & 3{,}674 & 13 & 4 & $34.1{\pm}7.7$ & $66.5{\pm}18.1$ & $44.37{\pm}16.0$ \\
\bottomrule
\end{tabular}
\end{table}

\subsection{Matched-scope results}

Our reruns approached published benchmarks for attempted speech,
typing, and handwriting. The remaining gaps were larger for inner
speech and voice synthesis, and the original handwriting model was
closer to its matched held-out benchmark before language modeling than
after it (see the handwriting comparison below).

\paragraph{Evaluation scope and session indexing.} Our initial Card et al.,
2024 decoding pipeline indexed the 45-session day-specific input layers
with the 41-session validation subset enumeration, silently routing
every validation session through an earlier day's affine layer (off by
one session in 2023 and up to four sessions in 2024--2025) and inflating
all-41-session pooled WER to 21.5--22.2\% (and 2025 follow-up sessions
to 67.3\%). Fixing \texttt{day\_idx} to the full 45-session index and aligning
study-day calendars against Table~S2 of Card et al., 2024 cut all-41-session
pooled WER by 13.5 percentage points to \textbf{7.02\%} with 5-gram +
OPT rescoring (8.70\% before rescoring, and \textbf{6.25\%} with a
bidirectional GRU encoder at 48k steps), matching the 6.70\%
Brain-to-Text~'25 baseline; across the 12 paired evaluation days in
their Table~S3 (sessions 2--17), our median WER is \textbf{5.58\%} against the
paper's online median of 4.49\% (with our WER lower on 4 of 12 days and
\textbf{3.48\%} on session~2, 2023-08-13, inside the paper's 9.8\%, CI
4.1--16.0). Conversely, the headline of Willett et al., 2023 is
Table~1's \emph{mean across all evaluation days} (23.8 under a 3-gram,
online); our per-day distribution matches it at the \emph{per-day
median}, 23.58---pooling sentences across days (the naive statistic)
inflates our number to 31.76 until the statistic is matched. The
statistic a paper reports is itself a reproduction parameter. Figure~\ref{fig:vocab}c uses one word-count convention for the released
online text outputs: within each recording day, sum substitutions,
deletions, and insertions and divide by the total reference words.
The bar is the equal-day mean and the whisker is one sample SD;
single-day results have no whisker. Text is case-folded and punctuation
is stripped within whitespace-delimited words, retaining apostrophes.
For Kunz et al., 2025, the selected inner-speech conditions give 22.9\%
(T12, one day), $43.3\pm23.8$\% (T15, two days), and 44.4\%
(T16, one day). Jude et al., 2026a gives $10.5\pm3.1$\% over T17's
six days and $3.5\pm0.6$\% over T18's seven days; Jude et al., 2026b
gives $66.5\pm18.1$\% over four days. Willett et al., 2021 gives
$23.9\pm8.4$\% over five days. Wairagkar et al., 2025 remains a
separate listener-transcription median (43.8\%), without a daily SD.

\paragraph{Session variability and release coverage.} After the
indexing fix, two-thirds of the apparent ``late-2025 collapse'' proved
to be an indexing artifact: the six 2025 follow-up sessions pool at
13.27\% WER (+OPT; 2023: 6.52\%, 2024: 5.15\%), with residual difficulty
concentrated in four specific sessions (2023-10-13, 2025-01-10,
2025-03-14, and 2025-03-30 at 20.2--26.0\% WER). Moreover, cross-checking
Table~S2 of Card et al., 2024 shows that the public release contains 32
of the paper's 84 in-window sessions ($\sim$38\%, plus 13 post-window
sessions in 2024--2025). Fine-tuning from a saturated 45-session
checkpoint on a hard day fails (near-zero gradient), whereas training a
44-session model that excludes 2025-01-10 and then cross-fitting on that
day's 106 training sentences drops its validation PER from 61.29\%
(zero-shot) to \textbf{26.03\%} (matching the 24.38\% full-soup baseline)---confirming
that fresh per-day supervision is worth $\sim$35 PER points on hard days.
The Fan et al., 2023 release shows the same coverage gap from the other
side: its two released evaluation arms (10 of 15 sessions, 20--40
sentences per session against the paper's average 58.3) rescore to
4.10\%/15.46\% WER---better than the paper's own 6.16\%/26.51\% precisely
because the shipped subset is easier than the full record.

\paragraph{Attempted and inner speech.} Our offline greedy-CTC runs
show higher phoneme error for inner speech within participants (T15
10.2$\to$59.7\% PER from overt copy to inner, T12 23.4$\to$61.4\%,
with T17 intended speech at 40.0\% and T16 inner speech at 61.5\%
landing in the same band, Fig.~\ref{fig:vocab}e), with the caveat
that our inner-speech arms train on hundreds of sentences against the
overt runs' full cumulative checkpoints. The remaining gap to
Kunz et al.'s online benchmark also reflects the language-model
substitution (the Willett 5-gram) and offline evaluation
(Table~\ref{tab:repro}). Section~\ref{sec:home_use} discusses the
within-participant online comparisons.

\paragraph{The original handwriting decoder.} We retrained the Willett
et al., 2021 decoder from scratch in PyTorch. The original pre-CTC
pipeline used labels from HMM forced alignment, a GRU frame
classifier, and a bigram language model. On the same held-out trials,
our raw CER was 4.71\% versus 2.78\% for the official checkpoint
(+1.93 percentage points, within the $\pm$2-point tolerance).
The bigram LM halves our
WER (19.70\% $\to$ 10.00\%), the same direction as the official
pipeline's larger gain (12.88\% $\to$ 3.64\%). The residual LM-tier gap
traces to one explainable artifact: the retrained model emits a
spurious character-onset spike at sentence onset, which the bigram LM
translates into a legal leading word; away from sentence onset, its
per-sentence errors are nearly character-identical to the official
model's.

\section{When language models help and what hardware reproduction requires}
\label{sec:feasibility}

Figure~\ref{fig:phase_memory} relates word errors after language
modeling to errors at the neural decoding stage (a,b) and shows the
memory needed to run the reproduced pipelines (c,d).

\begin{figure}[tp]
\centering
\includegraphics[width=\linewidth]{main_fig_3.pdf}
\captionsetup{singlelinecheck=false,font=footnotesize,skip=2pt}
\caption{\textbf{Reproduction feasibility: error-rate phase diagrams and
the memory wall.} (a) Session-level greedy phoneme error rate (PER)
versus post-LM WER across 77 attempted/intended-speech sessions in
three studies (Willett et al., 2023, T12; Card et al., 2024, T15;
Jude et al., 2026b, T17; open markers: sessions, stars: pooled;
dash-dotted: no-LM baseline $1-(1-\text{PER})^{3.4}$; dashed:
$\text{WER}=\text{PER}$). WER falls below PER in 45/57 sessions (79\%)
below the 25\% PER crossover and 4/20 (20\%) above it. (b) The same diagram for raw character error
rate (CER, no LM) versus post-LM WER across the four character-level
arms we reproduced (Willett et al., 2021, T5; Fan et al., 2023, T5;
Jude et al., 2026a, T17/T18; no-LM baseline exponent 5.1). (c, d) Peak
host RAM (c) and GPU VRAM (d) of the official stacks versus our
memory-mapped/KenLM single-GPU stack across the eight open studies
(dashed: 64\,GB RAM and 16\,GB VRAM consumer-PC ceilings): ours peaks
at 26.9\,GB RAM and 13\,GB VRAM versus 300\,GB and 80\,GB
official---every reproduction fits on one consumer PC (widest gap:
Wairagkar et al., 2025, 1.5 vs.\ 72\,GB VRAM).}
\label{fig:phase_memory}
\label{fig:pervswer}
\end{figure}

\subsection{Neural decoding errors and language-model gains}

Figure~\ref{fig:phase_memory}a plots, for every research session in our
three attempted- or intended-speech reproductions (77 sessions; Kunz
et al., 2025 omitted as inner speech), the session-level greedy PER
(acoustic stage only) against the session-level WER after the full
$n$-gram + rescoring pipeline. Two reference curves separate
granularity inflation from language-model correction. An English word
in these evaluation corpora averages $L = 3.4$ pure phonemes, so a
single phoneme error corrupts an entire word: under independent
phoneme errors a zero-LM decoder would follow the upper dash-dotted
curve $\text{WER}_{\text{no-LM}} = 1 - (1 - \text{PER})^{3.4}$
(greedy CTC decodes cluster errors within words, giving the
empirically milder $1 - (1 - \text{PER})^{2.3}$), while the
$45^\circ$ diagonal $\text{WER} = \text{PER}$ marks where word errors
after language modeling equal the raw phoneme error rate. All 77 sessions lie below the
modeled no-LM curve. Post-LM WER falls below raw PER at a crossover
near $\text{PER} \approx 25\%$ (logistic fit: 24.97\%; vertical gray
line). Below it, 45 of 57 sessions (79\%) reach
$\text{WER} < \text{PER}$: the Card et al., 2024 corpus drops from
10.3\% PER (21.8\% no-LM word error) to 7.02\% pooled WER, and
Willett et al., 2023 reaches a pooled 23.26\% WER from 23.4\% PER
(43.0\% without a language model). Above the crossover, 4 of 20 sessions
reach $\text{WER} < \text{PER}$. Large word errors can remain after
language modeling (Jude et al., 2026b pooled 40.0\% PER and 46.0\% WER;
Kunz et al., 2025 pooled WER 81--122\%).

Figure~\ref{fig:phase_memory}b compares raw character error rate (CER)
with post-LM WER for the reproduced handwriting and typing arms. Its
reference curve uses an average word length of 5.1 characters. Each
panel relates unit-level errors to word errors within its own set of
decoding tasks.

\subsection{RAM and GPU memory requirements}
\label{sec:memory}

Figure~\ref{fig:phase_memory}c,d compares peak host RAM and GPU
memory across the eight public studies. Our memory-mapped and KenLM
implementations peaked at 26.9\,GB RAM and 13\,GB VRAM across the
reproductions, compared with 300\,GB RAM and 80\,GB VRAM for the
official stacks. All eight reproductions fit within the 64\,GB RAM and
16\,GB VRAM limits of our PC.

Two 2026 preprints replaced a \texttt{wenet} fork of Kaldi/OpenFST
weighted finite-state transducers (WFST) and its $\sim$350\,GB
composed search graphs with KenLM \cite{heafield2011} behind Flashlight-Text, holding
WER constant (9.08\% $\to$ 9.03\% on the same 5-gram) while cutting RAM
to under 10\,GB, load time from $\sim$15 minutes to under 10 seconds,
and decoding latency by $3.11\times$ \cite{fogg2026, yoon2026}; the
rescoring LLM shrank from OPT-6.7B \cite{opt} to Qwen3.5-4B, and the
ensemble merger from a remote GPT-4.1-mini API to a locally served
Qwen3-4B \cite{fogg2026, yoon2026}. Two further preprints push the same
direction from outside the trial labs: LightBeam, a non-WFST CTC
decoder, cuts decoding RAM from $\sim$320\,GB to $\sim$10\,GB while
matching benchmark state of the art, by folding an LLM into beam search
through delayed fusion \cite{feghhi2026}; and an end-to-end decoder
with a cross-species pretrained encoder halves the prior end-to-end
WER on the open benchmarks, from 24.69\% to 10.22\% \cite{bit2026}.

\section{From laboratory demonstrations to independent home use}
\label{sec:maturity}

Figure~\ref{fig:clinical} follows three requirements for home use:
recording from appropriate sites over time (a--c), keeping setup and
recalibration manageable (d--f), and supporting the tasks users need
in everyday communication (g--j).

\begin{figure}[tp]
\centering
\includegraphics[width=\linewidth]{main_fig_4.pdf}
\captionsetup{singlelinecheck=false,font=footnotesize,skip=2pt}
\caption{\textbf{Clinical translation: the barriers between laboratory
demonstrations and a durable home device, grouped into three rows---
surgery-related (a--c), decoding technology (d--f), and patient
paradigm (g--j)---separated by dashed rules.}
(a) Non-speech cortical sites carry
no usable speech signal (paired bars: word-decoding accuracy from the
speech-capable site (solid) versus the non-speech site (hatched) of the
same participant): T12's ventral 6v arrays decode the 50-word task at
94\%, Broca's area~44 at 5.8\% (its movement/phoneme/sentence
accuracies: 11.5\%/9.6\%/$\sim$0\%) \cite{willett2023}; T16's i6v decodes
7-word attempted speech at 80.0\%, while the PEF/55b-border and 6d
hand-knob arrays stay at the $1/7$ word chance across all seven verbal
behaviors and were excluded from decoding \cite{kunz2025}; T17's 6v
reaches 51.1\% (95\% CI 48.0--54.2) on the 50-word vocabulary, dorsal
6d only 2.0\% (95\% CI 1.1--2.9) \cite{jude2026lis}. Ventral orofacial
6v/i6v decodes far above chance in all three cohorts; mid-precentral
55b only partially (T17) \cite{jude2026lis, deo2026}. (b) Chronic
Utah-array retention across 14 users ordered by implant duration
(0.8--7.4\,yr; hatched: $<1.5$\,yr, $n=7$; solid: $\ge 1.5$\,yr,
$n=7$) \cite{hahn2025}. (c) Channel counts across two craniotomies
(Table~\ref{tab:cohort}): B1 128 ch (4\,mm) in Moses et al., 2021
$\to$ 253 ch (3\,mm) in Brosler et al., 2026 (Bravo-1's
higher-density re-implant); T12 256 ch (L6v$+$L44, Willett et al.,
2023) $\to$ 512 ch (L6v$+$L44$+$R6v$+$R55b) after the 2025 bilateral
re-implant tested by Yoon et al., 2026
\cite{moses2021, willett2023, yoon2026, brosler2026}---when
yield decays beyond what (b) retains, the only current answer is more
surgery. (d) Cross-subject: single-user
vs.\ six-user-pretrained offline WER across the six BrainGate speech
users (digitized from Fogg et al., 2026 \cite{fogg2026}; relative WER
reductions, left to right: $-$31\% for T16, $-$57\% for T22, $-$50\%
for T12, $-$14\% for T15L, $-$30\% for T15S, $-$65\% for T17, and
$-$73\% for T21; open reanalysis T12, T15-S, T17; closed T15-L;
preprint T16, T21, T22; dashed line: 12.8\% Whisper large-v2
zero-shot WER averaged over 14 test sets (Radford et al., 2022,
Table~2) \cite{radford2022}; bars ordered
stroke (T16, T22) left and ALS (T12, T15L, T15S,
T17, T21) right of a central dashed divider). (e) New-user few-shot
calibration (Fogg et al., 2026, Fig.~5): first-day offline WER from
one session---single-user from scratch (light) vs.\ few-shot
adaptation of a frozen four-user pretrained decoder (dark; 95\%
CIs): at 90 sentences, T21 87.0\% vs.\ 11.2\% and T22 76.9\% vs.\
10.2\%; dashed line as in (d) \cite{fogg2026}. (f) Daily calibration
sentence requirements from 2021 to 2026, ordered
stroke (T22) / ALS (T12, T21, T17, T15) / SCI (T5, T18) around dashed
dividers, with T17 and T18 split from their shared 30-sentence
calibration condition
\cite{card2024, jude2026typing, willett2021, fan2023, fogg2026, card2026}.
(g) Copy-task benchmark (3.5\% WER at 49.7 \wpm{}, silent strategy)
versus self-generated free-use disagreement rate (46.7\%) during T15's
home use \cite{card2026} (home-use versus laboratory
corpora in Supplementary Fig.~S2c). (h) Online attempted vs.\ inner
speech at matched vocabularies
\cite{card2024, willett2023, kunz2025, nason2026}. (i) Online
within-participant comparison in the locked-in T17: attempted speech
$66.5\pm18.1$\% online WER over 3{,}674 released sentences
\cite{jude2026lis} vs.\ touch-typing $12.5\pm2.5$\% over 668 sentences
\cite{jude2026typing}. (j) Online output modality in T15:
brain-to-voice 43.8\% WER without a sentence-level LM
\cite{wairagkar2025} vs.\ brain-to-text 2.5\% with LM and LLM
\cite{card2024}; dashed line: offline Brain2voice 2.0 at 5.24\%
\cite{wairagkar2026}.}
\label{fig:clinical}
\end{figure}

\subsection{Implant placement and long-term recording}
\label{sec:implant_durability}

The within-participant comparisons in Fig.~\ref{fig:clinical}a show
that placement matters. Ventral orofacial sites (6v/i6v) decoded
attempted or intended speech substantially better than the comparison
sites in the same participants \cite{willett2023, jude2026lis, kunz2025}.

Across 14 BrainGate participants followed for up to 7.4 years, the fraction of
Utah-array electrodes recording clear neural spikes averaged 35.6\%,
declining from 41\% in months 1--3 to 34\% in the final months
(Fig.~\ref{fig:clinical}b) \cite{hahn2025}. Day-to-day stabilization can help compensate for
signal drift \cite{pun2024, karpowicz2025stab, wilson2025, hosman2023}, and
surface ECoG grids have supported a small command set over three
months without recalibration \cite{luo2023}.

A second operation increased channel counts in both B1 and T12
(Fig.~\ref{fig:clinical}c), through higher-density ECoG in B1 and
bilateral array coverage in T12
\cite{moses2021, willett2023, yoon2026, brosler2026}. These cases illustrate
the surgical burden of replacing or expanding implanted recording
hardware.

\subsection{Cross-user transfer and calibration}
\label{sec:calibration}

For new-user setup, early decoders relied on larger participant-specific
training corpora (\S\ref{sec:corpus_size}). Jude et al., 2026a
demonstrated first calibration for typing with 30 sentences
\cite{jude2026typing}. For speech, Fogg et al., 2026 used a multi-user
model to reach 10.2\% offline WER for a new user after 90 trials
($\sim$17 minutes; Fig.~\ref{fig:clinical}e) and 6.5\% after 200
\cite{fogg2026}. Figure~\ref{fig:clinical}d shows offline comparisons
between individual and pooled models.

Daily recalibration addresses changes in recordings from a user whose
decoder is already trained (Fig.~\ref{fig:clinical}f). Card et al., 2024
demonstrated rapid recalibration in 20 sentences (7.5 minutes)
\cite{card2024}. During T15's later home use, start-of-day speech
calibration used 0--20 optional sentences, with continued fine-tuning on
user-confirmed utterances and 9.2--16.9\,min/day of cursor calibration
\cite{card2026}.

Table~\ref{tab:crosssubject} places cross-user transfer for intracortical
decoders alongside related work on other interfaces. Peripheral interfaces set the scale
ceiling: a wristband sEMG decoder trained on 6{,}627 users works
zero-calibration on held-out users, while the fully open 108-user
emg2qwerty typing corpus quantifies how wide the cross-person domain
gap remains (55.4\% zero-shot CER against 6.95\% after $\sim$2.5\,h of
personal fine-tuning) \cite{kaifosh2025, sivakumar2024}. The
plausibility of pooling spike-resolution data across brains was built
in animals first---intracortical motor foundation models pretrained
across monkeys, human participants, and mice
\cite{azabou2023, ye2023, ye2025, zhang2025ned, fang2026}---but on motor
benchmarks whose metrics say nothing about communication, so
Table~\ref{tab:crosssubject} tracks only human studies. For
intracranial speech, cross-patient work on bulk field potentials has
so far been confined to overt speech and small vocabularies
\cite{makin2020, singh2025, spalding2026}; at action-potential
resolution in paralyzed participants, open-data reanalyses came first
\cite{boccato2026, bit2026}, and the six-user joint training of Fogg et
al., 2026 is the first in which every training participant benefited
and a held-out user crossed the high-accuracy line from 90 calibration
sentences \cite{fogg2026}. The bracketing domains also expose the
openness gradient: the largest pools---6{,}627 sEMG users, the
six-user intracortical corpus---are exactly the ones outsiders cannot
inspect, and the largest fully open human corpus remains a 108-user
sEMG typing set \cite{fogg2026, kaifosh2025, sivakumar2024}.

\begin{table}[tp]
\centering
\scriptsize
\setlength{\tabcolsep}{3.2pt}
\renewcommand{\arraystretch}{1.14}
\caption{\textbf{Cross-subject generalization across human
communication interfaces: pool, new-user adaptation, headline outcome,
and data release.}
Groups run from the peripheral engineering reference to intracortical
communication BCIs in paralyzed participants (whose IDs refer to
Table~\ref{tab:cohort}). \textbf{N} = users or animals in the training
pool (n/r = not reported); \textbf{Population} = who they are; \textbf{Pool} = training
corpus; \textbf{New-user data} = what a held-out
user must supply; \textbf{Zero-shot} = usable decoding for a new user
with none of their own data (yes/no; emg2qwerty's zero-shot CER of
55.4\% exists but is not usable, hence no); \textbf{Few-shot} = the
new user contributes some of their own labeled data---more than none,
at most one recording hour ($0 < \text{data} \le 1$\,h)---converted
at each task's own pace (yes = a $\le$1\,h adaptation of a held-out
individual is demonstrated; unreported amounts count as no);
\textbf{Headline outcome}
= one metric per study---WER for large-vocabulary speech, CER for
typing/handwriting, PER for phonemes, accuracy for closed
sets---modality-specific and \emph{not} comparable across rows,
offline unless noted; \textbf{Data} = release status of the neural
data itself (open / partial: some data or code public / closed).}
\label{tab:crosssubject}
\begin{tabular}{@{}>{\raggedright\arraybackslash}p{2.45cm}>{\centering\arraybackslash}p{0.65cm}>{\raggedright\arraybackslash}p{1.50cm}>{\raggedright\arraybackslash}p{1.45cm}>{\raggedright\arraybackslash}p{1.70cm}>{\centering\arraybackslash}p{0.75cm}>{\centering\arraybackslash}p{0.75cm}>{\raggedright\arraybackslash}p{3.60cm}>{\centering\arraybackslash}p{0.75cm}@{}}
\toprule
Study & $N$ & Population & Pool & New-user data & \shortstack[c]{Zero-\\shot} & \shortstack[c]{Few-\\shot} & Headline outcome & Data \\
\midrule
\multicolumn{9}{@{}l}{\emph{Automatic speech recognition (engineering reference)}} \\
Radford et al., 2022; Baevski et al., 2020 \cite{radford2022, baevski2020} & n/r & speakers (audio) & internet-scale & none & yes & no & speaker-independent transcription (deployed) & partial \\
\midrule
\multicolumn{9}{@{}l}{\emph{Surface EMG, able-bodied users}} \\
Sivakumar et al., 2024 \cite{sivakumar2024} & 108 & able-bodied & 346\,h & fine-tune, $\sim$2.5\,h & no & no & CER 55.4\%$\to$6.95\% (offline) & open \\
Kaifosh \& Reardon, 2025 \cite{kaifosh2025} & 6{,}627 & able-bodied & proprietary corpus & none & yes & no & 20.9 \wpm{} air-writing (closed-loop) & partial \\
\midrule
\multicolumn{9}{@{}l}{\emph{Bulk-field intracranial, clinical cohorts (overt speech)}} \\
Makin et al., 2020 \cite{makin2020} & 4 & epilepsy & 30--50 sentences & $\le$40\,min & no & yes & WER 3\% (offline, closed set) & partial \\
Singh et al., 2025 \cite{singh2025} & 25 & epilepsy & tongue twisters & $\sim$1\,h & no & yes & PER 57\%$\to$45\% (offline) & closed \\
Spalding et al., 2026 \cite{spalding2026} & 8 & surgery & 52 non-words & $\sim$5\% of data & no & yes & acc.\ 0.31 $>$ 0.24 within-patient (offline) & partial \\
\midrule
\multicolumn{9}{@{}l}{\emph{Intracortical communication BCIs, paralyzed participants}} \\
Boccato et al., 2026 \cite{boccato2026} & 2 & BCI users & T12+T15 corpora & linear fit, 12\,min--3\,h & no & yes & matches within-subject (offline) & open \\
Y.~Zhang et al., 2026 (BIT) \cite{bit2026} & 2 & ALS users & cross-species & thousands of sentences & no & no & WER 24.69\%$\to$10.22\% (offline) & open \\
Levin et al., 2026 \cite{levin2026} & 5 & BCI users & 5 public corpora & $<$200 sentences & no & yes & PER 23--40\% with transfer (offline, no LM) & partial \\
Fogg et al., 2026 \cite{fogg2026} & 6 & BCI users & joint 6-user corpora & 90 sentences & no & yes & $-$51.1\% rel.\ WER (offline) & partial \\
\bottomrule
\end{tabular}
\end{table}

\subsection{Communication in everyday use}
\label{sec:home_use}

Of the participants in Table~\ref{tab:cohort}, \textbf{one} has
documented independent use at home. BrainGate participant T15 adopted a
combined speech-and-cursor neuroprosthesis as his \emph{primary}
communication device and computer interface in daily life over 22.6
months, logging more than 3{,}800 hours of use across 444 of 653 days
(averaging 9.5\,h/day once his care partners could connect and start the
system without research staff) and producing 183{,}060 self-rated
sentences \cite{card2026}. During structured sentence-copy tests, his
system achieved $>$99\% accuracy when vocalizing (30.6 \wpm{}) and
96.5\% accuracy when speaking silently (49.7 \wpm{}, 3.5\% WER).
In unscripted daily conversation, 46.7\% of decoded utterances required
some user edit or differed from a post-hoc transcript. The two
conditions in Fig.~\ref{fig:clinical}g use different measures: word
error for sentence copying and utterance-level disagreement for free
use \cite{card2026}. T5's 403-day
plug-and-play handwriting demonstration took place during monthly,
researcher-supervised laboratory sessions rather than independent home
use \cite{fan2023}, and Neuralink's statement that P7 uses decoded
ten-finger typing for everyday computer control \cite{neuralink2026}
is not yet accompanied by a quantitative usage log. Every other result
in the literature is a supervised laboratory demonstration.

\paragraph{Attempted and inner speech.} Figure~\ref{fig:clinical}h
pairs online inner-speech WER with an articulation-engaging benchmark
for each participant. On the same 50-word sentence set, T12 reached
9.1\% WER when trying to vocalize \cite{willett2023} and 23.7\% with
inner speech \cite{kunz2025}. At matched vocabulary sizes, T15
reached 4.9$\pm$3.3\% with attempted speech \cite{card2024} and
51.3$\pm$33.1\% with inner speech across two 125k sessions. T16
reached 19.6\% with mimed speech \cite{nason2026} and 46.6\% with
inner speech \cite{kunz2025}. Kunz et al. also interleaved
attempted-vocalized and inner-speech conditions within experimental
blocks \cite{kunz2025}.

\paragraph{Comparing the three modalities for clinical readiness.}
\emph{Speech decoding} has been demonstrated in the largest number of
participants ($n=16$). Performance varies across users and recording
conditions (\S\ref{sec:defs}). Multi-user pretraining brought the
locked-in participant T17 to 23.1\% WER \cite{fogg2026}. In a conference
abstract, the Neuralink VOICE trial reported 5.8--6.6\% WER at
54--75 \wpm{} with fully implanted wireless arrays
\cite{pouratian2026}. \emph{Touch-typing decoding} reached
1.6\% WER after 30 calibration sentences and 81\% of able-bodied
smartphone typing speed in a task with 30 discrete output keys
\cite{jude2026typing}. The published study includes two participants,
with a third reported by a company. Typing speed was lower in the
participant with ALS (T17, 47
\cpm{}) than in the participant with spinal cord injury (T18, 110
\cpm{}), suggesting that typing accuracy depends on how well finger
representations are preserved in the motor cortex. \emph{Handwriting
decoding} achieved 90 \cpm{} in T5 \cite{willett2021}, but its original
two-stage character-recognition algorithm has now been replaced by the
same sequence-decoding architecture used for speech and typing
\cite{levin2026}. Across these demonstrations, modality choice depends
on what movements the participant can attempt and how the interface
fits their communication needs.

\subsection{Evidence for regulatory approval}

Every study reviewed
here was conducted under an investigational clinical-trial exemption
(such as FDA Investigational Device Exemption NCT00912041 for
BrainGate2, whose interim safety data are published \cite{rubin2023});
no implanted communication neuroprosthesis has received commercial
medical-device approval anywhere in the world. Five practical barriers
still separate current feasibility studies from a regulatable medical
product:
\begin{enumerate}[leftmargin=1.5em,itemsep=2pt]
\item \emph{Multicenter clinical trials beyond single-lab case series.}
All published evidence comes from small single-center cohorts, and even
the largest cross-patient speech decoder has been trained on only six
human brains \cite{fogg2026} (well below the tens of subjects used in
animal neural-foundation models \cite{ye2025}). Regulatory approval will
require prospectively powered, multi-site trials enrolling tens of
patients stratified by disease cause (ALS vs.\ brainstem stroke vs.\
spinal cord injury).
\item \emph{Long-term implant durability standards.} Electrode
retention and reimplantation results (\S\ref{sec:implant_durability})
motivate a quantitative standard for how long a patient can rely on the
device before surgical replacement is needed.
\item \emph{Minimal daily calibration and caregiver independence.}
Asking a paralyzed patient to spend hours recalibrating a decoder every
morning is incompatible with home care. As shown in
\S\ref{sec:calibration}, daily calibration has dropped to 20--30
sentences (or optional zero-calibration days with background updating on
user-confirmed sentences during T15's home use \cite{card2026}),
demonstrating that low-burden daily operation is achievable.
\item \emph{Fully implanted wireless hardware and compact computing.}
Moving out of the laboratory requires replacing percutaneous head
pedestals and server racks with wireless implants \cite{pouratian2026, simeral2021}
and portable processors. Both the field's recent switch to lightweight
language models \cite{fogg2026, yoon2026} and our single-GPU
reproductions show that the entire decoding software stack now fits on a
single consumer graphics card.
\item \emph{Safe handling of decoding failures in daily life.} When a
user switches between vocalized and silent speech strategies or skips
calibration across several days, decoder accuracy can drop sharply---a
failure mode seen both in clinical trials and in our reproductions. A
commercial medical device will need automated quality monitoring and
graceful fallback modes, potentially aided by decoding neural
error-detection signals directly from the user's motor cortex
\cite{hou2025}.
\end{enumerate}
Coordinated multicenter trials would allow these requirements to be
evaluated across diagnoses, implant configurations, and home-use
settings.

\section{Designing speech neuroprostheses for Mandarin}
\label{sec:chinese}

Nearly every communication neuroprosthesis to date has been built for
English, yet roughly 7{,}000 spoken languages worldwide remain
unaddressed, and China alone has the world's largest patient population
living with severe speech paralysis. Extending communication
neuroprostheses to Mandarin Chinese requires rethinking both the
clinical target population and the linguistic design of the decoder
(Fig.~\ref{fig:mandarin}a--e).

\begin{figure}[tp]
\centering
\includegraphics[width=\linewidth]{main_fig_5.pdf}
\captionsetup{singlelinecheck=false,font=footnotesize,skip=2pt}
\caption{\textbf{Future of Chinese: our insights toward Mandarin speech
neuroprostheses.}
(a) Annual
survivors in China: brainstem stroke vs.\ ALS
($\sim$240{,}000--280{,}000/yr chronic brainstem stroke survivors vs.\
$\sim$20{,}000--30{,}000/yr ALS survivors)---ALS progresses within
2--5 years while brainstem stroke is non-progressive with 10--20+ year
survival, so stroke dominates the addressable population and sets a
10-year-plus chronic-durability requirement on the implant. (b) Phoneme
count: English's 39 ARPAbet phonemes (24
consonants $+$ 15 vowels; 41 CTC classes with silence and blank) vs.\
Mandarin's 57 phonological initials $+$ finals (21 $+$ 36, derived
from the 39-final inventory of the standard \emph{Xiandai Hanyu}
textbook \cite{huangliao2017} by dropping marginal \^{e}, \emph{er},
and \emph{ueng}; the 47
literal Pinyin letters, 23 $+$ 24, conflate the three \emph{i}-vowels
and [u]/[y])---the decoder's output inventory should follow the
phonological count, not the orthographic one. (c) Native
vocabulary count: 3{,}500 Tier-1 common characters (Chinese standard,
$>$99\% character token coverage) and the full 8{,}105-character 2013
official standard (Table of General Standard Chinese Characters; three
tiers $3{,}500 + 3{,}000 + 1{,}605$; $\geq$99.9\% token coverage)
against the three English tiers of Brysbaert et al., 2016
\cite{brysbaert2016}---11{,}100 word families, 42{,}000 lemmas,
71{,}400 inflected word forms ($1.7\times$ the lemma count);
character-level composition compresses the lexicon a Chinese decoder
must map. (d) The
Mandarin Pinyin trap (midsagittal vocal-tract kinematics): one letter
\emph{i} labels three opposite gestures---\emph{ji} (tongue body raised
toward the hard palate, lips spread, $F_2 \approx 2{,}200$--$2{,}500$
Hz), \emph{si} (apical vowel, tongue tip flat against the alveolar
ridge, $F_2 < 1{,}200$ Hz, the lowest of the three), and \emph{shi}
(retroflex vowel, tongue tip curled up and back, $F_2 \approx
1{,}200$--$1{,}600$ Hz); literal-Pinyin labels would force three
distinct motor plans into one class, so a Mandarin speech BCI requires
articulatory-faithful initial/final variant labels. (e) Mandarin 4
lexical tones: Standard Mandarin has four lexical
tones (other spoken Chinese varieties such as Cantonese or Min
differ); curves follow measured time-normalized F0 contours
\cite{tupper2020, shiwang2006}---Tone 2 dips before rising, the Tone 3
trough sits mid-syllable with only a partial final rise, and Tone 4
holds its high onset before an accelerating fall---and this four-tone
axis is already decoded in Mandarin speech
BCIs---93\% tone accuracy (chance 25\%) from onset--syllable--tone
decomposition over five ECoG participants \cite{zhangd2024}, and a
real-time dual-stream decoder of 394 syllables $\times$ 4 tones on a
flexible 256-channel micro-ECoG array at 71.2\% median syllable
accuracy \cite{qian2025}.}
\label{fig:mandarin}
\end{figure}

\subsection{Clinical population and implant durability}

In both China and
the United States, brainstem stroke causes roughly ten times as many new
cases of severe speech paralysis each year as ALS ($\sim$300{,}000--350{,}000
vs.\ $\sim$20{,}000--30{,}000 annual cases in China;
$\sim$70{,}000--80{,}000 vs.\ $\sim$5{,}000--6{,}000 in the US;
Fig.~\ref{fig:mandarin}a). Moreover, whereas ALS is a progressive disease
with a typical survival of 2--5 years, brainstem stroke is
non-progressive and survivors often live for 10--20 years or longer.
Brainstem stroke is therefore the largest clinical indication for a
long-term communication neuroprosthesis, making decade-scale hardware
durability essential. Patient demographics also shape which movement
modality works best: while 80.72\% of people in China speak standard
Mandarin (Putonghua) \cite{moe2022lang} and 97.3\% are literate (and
thus learned handwriting in school), we estimate that only about 57.1\%
type Chinese using a Pinyin keyboard (calculated from national internet
adoption multiplied by the share of Pinyin input-method users
\cite{cnnic2025, iimedia2022}, as many older or rural users rely on voice
messaging or handwriting input on touchscreens; Supplementary Fig.~S10).
For older stroke survivors who never touch-typed on a QWERTY keyboard,
attempted speech and handwriting are far more natural targets than
touch-typing.

\subsection{Articulation, lexical tone, and character output}

\paragraph{Vocabulary and syllable structure.} At the language-model
stage, Chinese actually simplifies vocabulary coverage compared with
English (Fig.~\ref{fig:mandarin}c). Because Chinese words are composed
from a compact set of characters, the 3{,}500 Tier-1 standard characters
already cover $>$99\% of everyday written text, and the full
8{,}105-character national standard (Table of General Standard Chinese
Characters) covers $\geq$99.9\%---an order of magnitude smaller than the
71{,}400 inflected word forms known by an average English speaker
\cite{brysbaert2016}. At the acoustic-decoding stage, Mandarin's regular
syllable structure ($\sim$400 base syllables formed by combining 21
initial consonants with 36 final vowel/nasal rhymes
\cite{huangliao2017}, compared with English's 39 ARPAbet phonemes;
Fig.~\ref{fig:mandarin}b) naturally favors decoding initials, finals,
and lexical tones. Standard Mandarin uses four lexical tones (pitch
contours that distinguish word meaning, whereas other Chinese varieties
such as Cantonese use more; Fig.~\ref{fig:mandarin}e and Supplementary
Fig.~S11). Early studies in presurgical epilepsy patients with intact
speech have already shown that cortical surface recordings can decode
Mandarin tones at 93\% accuracy (chance 25\%) \cite{zhangd2024} and
stream 394 syllables $\times$ 4 tones in real time from a 256-channel
flexible ECoG grid \cite{qian2025}, confirming that pitch-contour
commands are represented in human sensorimotor cortex.

\paragraph{Articulatory labels for Pinyin finals.} A critical
pitfall for Mandarin speech BCIs is confusing romanized Pinyin spelling
with actual vocal-tract movements (Fig.~\ref{fig:mandarin}d). In Pinyin,
the single letter \emph{i} is used to write the vowel in three
completely different tongue postures: \emph{ji} (tongue body raised high
against the hard palate with spread lips), \emph{si} (tongue tip pressed
flat behind the front teeth), and \emph{shi} (tongue tip curled upward
and backward). Because the speech motor cortex encodes the physical
movement of the tongue and lips rather than how a word is spelled,
training a neural decoder to predict literal Pinyin letters would force
three distinct brain signals into a single label. A Mandarin speech
neuroprosthesis must therefore label its output units by true phonetic
articulation (21 initials $+$ 36 articulatory finals $+$ 4 tones) rather
than orthographic Pinyin letters.

\section{Limitations of the evidence}
\label{sec:limitations}

Four limitations should be kept in mind when interpreting our findings:
\begin{enumerate}[leftmargin=1.5em,itemsep=2pt]
\item \emph{Small, single-center clinical cohorts.} Nearly all patient
data in this field come from a few academic medical centers (primarily
the BrainGate consortium and UCSF). Because we synthesize published
studies without adding new clinical participants of our own, our
maturity assessment inherits the enrollment criteria of those trials
(such as selecting motivated participants with strong family support and
favorable cortical anatomy).
\item \emph{Reproduction coverage is limited to open intracortical
data.} We could independently reproduce only the studies that released
public datasets. Because the UCSF ECoG studies, the T15 home-use
recordings \cite{card2026}, and the six-patient corpus of Fogg et al.,
2026 \cite{fogg2026} are not publicly available, our reproduction audit
tests the open intracortical literature rather than the closed surface-ECoG
literature. Furthermore, where our inner-speech reproduction fell short
of the published online benchmark (Kunz et al., 2025 \cite{kunz2025};
Table~\ref{tab:repro}), the gap reflects a combination of smaller
training sets, language-model substitution, and offline evaluation.
\item \emph{Public datasets contain processed features, not raw brain
signals.} To protect patient privacy and reduce file size, every public
dataset releases pre-computed neural features (20\,ms bins of spike
counts and high-frequency band power) rather than raw broadband voltage
waveforms. Our reproductions therefore validate the machine-learning
decoding pipelines, but cannot independently audit the upstream signal
filtering or spike detection.
\item \emph{Cross-study differences in test protocols.} Even when we
match evaluation statistics (such as comparing per-day medians in
Willett et al., 2023 \cite{willett2023} or single-session benchmarks in
Card et al., 2024 \cite{card2024}), different studies still use
different sentence prompts, vocabulary sizes, and online versus offline
scoring rules. Finally, several 2026 advances reviewed here are
currently preprints on bioRxiv or medRxiv, and their exact numbers may
be refined during peer review.
\end{enumerate}

\section{Conclusion}

\begin{enumerate}[leftmargin=1.5em,itemsep=2pt]
\item \emph{Algorithmic convergence:} Speech, touch-typing, and
handwriting neuroprostheses have converged on one decoding pipeline
adapted from automatic speech recognition (binned motor-cortex activity,
a recurrent or Transformer neural network trained with connectionist
temporal classification, $n$-gram beam search, and large-language-model
rescoring), while replacing memory-heavy speech-recognition toolkits
with lightweight decoders in 2026.
\item \emph{Two-track data openness:} Intracortical microelectrode
studies since 2023 have broadly shared processed neural features on
public repositories (Dryad and Zenodo), whereas surface ECoG studies and
home-use recordings remain restricted, and no study has released raw
voltage waveforms.
\item \emph{Independent reproducibility on everyday hardware:} All eight
publicly available datasets can be retrained and evaluated from scratch
on a single 16\,GB consumer GPU, matching or approaching published
benchmarks on attempted speech, typing, and handwriting while confirming
that language models amplify usable neural signals rather than rescuing
noisy ones.
\item \emph{Calibration and transfer:} Early decoders trained on
thousands of trials from an individual participant. Recent studies
demonstrate daily recalibration with tens of sentences and new-user
speech adaptation with 90 trials after multi-user pretraining.
\item \emph{Clinical reality and path to approval:} Across five years
worldwide, 20 people with paralysis have used a communication
neuroprosthesis in real time (including two in company demonstrations),
and only one patient (T15) has used a system independently at home as a
primary communication device. Moving from laboratory feasibility to
regulatory approval now depends on multicenter clinical trials and
long-term implant durability standards.
\item \emph{Mandarin decoder design:} The clinical population in China
motivates decade-scale implant durability. The decoder must represent
articulatory distinctions and lexical tone and generate character-based
text.
\end{enumerate}

\clearpage
\section*{Data and code availability}

{\setlength{\leftskip}{1.2pt}\setlength{\rightskip}{1.2pt}\relax
All neural datasets analyzed in this review are publicly available from
the Dryad, Zenodo, and OSF repositories cited in \S\ref{sec:openness}
\cite{card2024, willett2023, wairagkar2025, jude2026lis, jude2026typing, willett2021, fan2023, kunz2025, btot}. Our reproduction code and per-trial
decode outputs are available at
\url{https://github.com/Physics-Lee/communication-neuroprostheses-reproduction},
and the milestone timeline database and rendering code for
Fig.~\ref{fig:timeline} are available at
\url{https://github.com/Physics-Lee/communication-neuroprostheses-timeline}.\par}

\section*{Acknowledgments}

We thank the twenty clinical trial participants, their families, and
their care partners, as well as the clinical research teams that
publicly shared their datasets.

\section*{Abbreviations}

\noindent
\begin{tabular}{@{}>{\raggedright\arraybackslash}p{1.6cm}>{\raggedright\arraybackslash}p{5.6cm}>{\raggedright\arraybackslash}p{1.6cm}>{\raggedright\arraybackslash}p{5.6cm}@{}}
AAC & augmentative and alternative communication & PER & phoneme error rate \\
ALS & amyotrophic lateral sclerosis & RNN-T & recurrent neural-network transducer \\
BCI & brain--computer interface & SCI & spinal cord injury \\
cpm & characters per minute & sEEG & stereoelectroencephalography \\
CTC & connectionist temporal classification & TLG & token--lexicon--grammar (WFST composition) \\
ECoG & electrocorticography & vPCG & ventral precentral gyrus (speech motor cortex) \\
GRU & gated recurrent unit & WER & word error rate \\
IDE & investigational device exemption & WFST & weighted finite-state transducer \\
LIS & locked-in syndrome & wpm & words per minute \\
LLM & large language model & & \\
\end{tabular}

\bibliographystyle{naturemag}%
\nocite{silva2025,candrea2025,barrow2025,neuralink2025e,seide2011,hinton2012,hannun2014,graves2012,he2019,chan2016,gulati2020,hsu2021,chu2023,rubenstein2023,tang2023,amodei2016,jiemian2020,cyd2013}%
\begingroup\footnotesize\sloppy\emergencystretch=1.5em
\microtypesetup{protrusion=false}%
\bibliography{refs}
\endgroup


\end{document}
